\section{Attack-Screening Addendum}
\label{app:attack-screening}

This appendix records the first validation step after the initial construction: a conservative attack-screening layer for parameter rejection.  The estimates here are not security proofs and do not certify concrete SABLE parameters.  Their purpose is to expose weak settings before implementation or benchmarking.  In particular, any final parameter set must still be analyzed by specialists in sparse-LPN, q-ary LPN, information-set decoding, and large-sample attacks.

\subsection{Screened attack surfaces}

The validation repository models two public noisy-linear-equation surfaces.  First, the expansion key exposes approximately
\[
  m_{\exp}=(n+1)^2
\]
sparse-LPN-like rows related to the GSW expansion layer.  Second, the compaction key exposes
\[
  m_{\mathrm{cmp}}=(n+1)m_c
\]
q-ary LPN-like samples for the code/LPN compactor.  These are only proxy surfaces; the exact joint distribution of the GSW rows and encrypted secret-key coordinates must still be handled in the formal proof and in cryptanalysis.

The repository now implements three first-pass attack screens.

\paragraph{No-clean-subset screen.}
A random set of $n$ noisy equations is clean with probability roughly $(1-\eta)^n$.  The corresponding infinite-sample work proxy is
\[
  \log_2 W_{\mathrm{clean}}
  \approx -n\log_2(1-\eta)+\omega\log_2 n,
\]
where the last term represents dense linear algebra.  This screen is deliberately simple and is especially important when correctness pushes $\eta$ to be very small.

\paragraph{Prange/information-set screen.}
Viewing the public samples as a noisy random linear codeword, a Prange-style information-set decoder succeeds when it chooses an error-free information set.  With $m$ samples and expected error weight $t=\eta m$, the proxy work is
\[
  \log_2 W_{\mathrm{Prange}}
  \approx
  \log_2 {m\choose n} - \log_2 {m-t\choose n} + \omega\log_2 n.
\]
This is the most basic ISD screen, following the information-set viewpoint of Prange and later code-decoding attacks~\cite{Prange1962,Stern1989,BerlekampMcEliecevanTilborg1978}.

\paragraph{BKW-style bias screen.}
For q-ary symmetric noise, the estimator uses the character-bias proxy
\[
  \beta_q = \left|1-\frac{q\eta}{q-1}\right|.
\]
After approximately $2^a$ combined samples, the proxy bias is $\beta_q^{2^a}$ and a final distinguisher wants on the order of $\beta_q^{-2^{a+1}}$ samples.  The code minimizes a simple table-building and distinguishing proxy over the number of BKW levels.  This follows the role of BKW-family algorithms in LPN cryptanalysis~\cite{BlumKalaiWasserman2003,LevieilFouque2006,BogosEtAl2015}.  It is not a specialized sparse-q-ary-LPN estimator.

\subsection{Early finding: correctness/security tension}

The first estimator run rejects the earlier \texttt{prototype\_medium} preset.  The reason is not a software bug: the preset used very low noise for correctness.  For $q=65537$, $n=512$, $k=3$, and $\eta=2^{-20}$, the expansion-key surface has expected total errors below one.  The no-clean-subset and Prange-style proxies therefore fall far below a 128-bit target.  The compaction-key surface has the same problem.

The repository records the following conservative feasibility grid.  Here $\eta$ is set to the largest value allowed by the conservative requirement that the evaluated GSW error be at most $0.10$ for one additive term.  The column $n_{\min}$ is the smallest dimension making the no-clean-subset proxy reach 128 bits.  The final column is the sparse expansion-key entry count $(n+1)^2(k+1)$.

\begin{center}
\begin{tabular}{rrrrr}
\toprule
$k$ & depth & $\eta$ ceiling & $n_{\min}$ & expansion entries \\
\midrule
1 & 1 & $6.67\cdot 10^{-3}$ & $9{,}172$ & $1.68\cdot 10^8$ \\
2 & 1 & $2.50\cdot 10^{-3}$ & $23{,}388$ & $1.64\cdot 10^9$ \\
3 & 1 & $1.18\cdot 10^{-3}$ & $47{,}905$ & $9.18\cdot 10^9$ \\
1 & 2 & $3.92\cdot 10^{-4}$ & $135{,}802$ & $3.69\cdot 10^{10}$ \\
2 & 2 & $3.05\cdot 10^{-5}$ & $1{,}510{,}083$ & $6.84\cdot 10^{12}$ \\
3 & 2 & $4.58\cdot 10^{-6}$ & $8{,}895{,}213$ & $3.16\cdot 10^{14}$ \\
\bottomrule
\end{tabular}
\end{center}

The table suggests that the current construction is most plausible for depth-one arithmetic, i.e., degree-two computations, unless a new refresh, coding, or compaction technique improves the correctness/security tradeoff.  Depth two is not ruled out mathematically, but the naive expansion-key size becomes very large under this conservative screen.

\subsection{Design consequence}

The next research problem is therefore sharper than the original construction problem:

\begin{quote}
Can SABLE reduce its dependence on extremely low LPN noise while preserving low-depth homomorphic correctness?
\end{quote}

Promising directions are: using a stronger q-ary code rather than repetition in the compactor; reducing the public sample exposure through seeded or punctured evaluation keys; adding a lightweight code-based refresh for the GSW layer; deriving sharper circuit-aware error bounds instead of worst-case balanced-product bounds; and testing whether the sparse-LPN assumption used by Corrigan-Gibbs, Henzinger, Kalai, and Vaikuntanathan~\cite{CorriganGibbsHenzingerKalaiVaikuntanathan2024} admits concrete parameter regimes compatible with code/LPN compaction.
